\section{Test og verifikation}

\emph{Test og verifikation} vurderer, om de centrale tekniske påstande om TræningsMester er dokumenteret. Målet er teknisk plausibilitet og tydelig afgrænsning af evidensniveauet.

Verifikationen læses som et afgrænset teknisk evidensspor. Den kan vise, at centrale dele kan bygges, at udvalgt domænelogik er testet, og at brugerflader og arkitektur er dokumenteret. Spørgsmål om motivation, fastholdelse og brugeraccept kræver senere bruger- og effektstudier.

\subsection{Build, unit tests og centrale flows}

Tabel~\ref{tab:tm-test-traceability} viser, at iOS-klient, watchOS-target og Live Activity-target kan bygges, og at central forretningslogik er \emph{unit-testet}. Det understøtter, at den udviklede app har implementeret domænelogik for tracker, completion, cycle runtime og Live Activity-projection.

\emph{Buildstatus} er afgrænset til den gennemgåede tilstand. Den dokumenterer kompilerbare targets, mens uafhængig \emph{CI-proces}, komplet installationsflow og stabil runtime på tværs af enheder kræver særskilt test. Unit tests dokumenterer afgrænsede logikker. Forståelighed og backend-/platformscenarier kræver andre testformer.

Unit tests passer især til de dele, hvor forventet adfærd kan beskrives præcist. Det gælder payloads, trackerstatus, completion, cycle runtime og Live Activity-projection. Spørgsmål om forståelighed, tempo og oplevet friktion kræver brugerafprøvning.

\subsection{Integrationstest og sikkerhedsbegrænsninger}

\emph{Integrationstests} er mere begrænsede. Tabel~\ref{tab:tm-test-traceability} angiver, at Supabase anon read path er delvist bestået, mens admin- og ikke-adminscenarier kræver yderligere dækning. Backend- og rolleverifikation står derfor som et vigtigt næste testspor.

Det er især vigtigt for sikkerhedsmodellen. Figur~\ref{fig:tm-security-boundary-summary} og Tabel~\ref{tab:tm-rls-matrix} viser en arkitektur, hvor klient, RLS og Edge Functions har forskellige roller. OWASP MASVS og Supabase-dokumentationen understøtter princippet, mens sikkerhedsbeviset kræver mere målrettet test af admin-, træner-, tværbruger- og Edge Function-scenarier \cite{owasp2024_masvs,supabase_rls_docs,supabase_securing_api_docs}.

\subsection{Visuel dokumentation og åbne mangler}

Skærmbilleder og diagrammer dokumenterer, at centrale brugerflader og arkitekturspor findes. Figur~\ref{fig:tm-ios-home-ready} viser Home som næste-handling-flade, Figur~\ref{fig:tm-ios-tracker-active} viser aktiv tracker, og Figur~\ref{fig:tm-system-context} viser systemkonteksten. Det styrker forbindelsen mellem krav, implementering og synlige artefakter.

Rest-risici handler især om reproducerbarhed, backendintegration, rollebaseret adgang og runtime for watchOS og Live Activity. Åbne mangler afgrænser, hvad der er dokumenteret. Her er det teknisk plausibilitet for centrale flows, mens sikkerhedscertificering, penetrationstest og driftssikker produktionsmodenhed kræver yderligere arbejde.

Den vigtigste konklusion fra verifikationen er dobbelt. Kerneflowet har et dokumenteret teknisk grundlag, men flere integrations- og sikkerhedsscenarier skal testes, før appen kan vurderes som moden til bred drift.

\subsection{Verifikationsvurdering}

Belægget bliver stærkere, når påstandene forbindes med konkrete artefakter som build, unit tests, skærmbilleder, diagrammer og sporbarhedstabeller. Dermed står TræningsMester som mere end en idé. Centrale dele af designet er omsat til en fungerende teknisk struktur.

Samtidig skal verifikationen bruges præcist. De stærkeste konklusioner handler om implementeringsplausibilitet, arkitekturvalg og kendte begrænsninger. Langtidseffekt, fuld sikkerhed og stabil runtime på alle enheder kræver bredere evaluering. Verifikationen er derfor et kontrolleret evidensspor.

\tmdelkonklusion{Verifikationens grænse}{Testene gør det rimeligt at tale om en fungerende teknisk struktur. Færdig drift, fuld sikkerhed og dokumenteret brugeradfærd kræver bredere test og evaluering. Den skelnen er central for hele konklusionen.}
